\section{Global row and source bounds on the actual sector}
\label{sec:actual-global-row-prices}

The palette comparison needs estimates on the whole selected row matrix,
including the source created by frozen links. A small covariance error
on a posed clean panel does not supply these estimates. We establish the
required global bounds using the actual sector precision and the selected
background rows.

Fix the gauge group, a one-step patch with block length $L$, its kernel,
and its funded tower and path prices. Assume the patch satisfies the
predictor branch conditions. All constants below are chosen from these
fixed data. Put
\[
 B=\lceil(\log x)^2\rceil,\qquad t=\operatorname{fiery}(x,L),
 \qquad R=\lfloor B/100000\rfloor.
\]
Here $t$ is the fine small-field threshold of
Section~\ref{sec:seed-guides}, while $\operatorname{temp}(x)$ is the
coarse threshold; eventually $\operatorname{temp}(x)\le2t$.
The coarse period satisfies $c\ge100(1+B)$. Use arbitrary stores,
forest, predictor tags and coarse field $V$ with the stores' flags equal
to the actual coarse flags at threshold $\operatorname{temp}(x)$.
No assertion about their fine costs or other noise entries is required.
The large-$x$ threshold is uniform over all these choices and all funded
towers having the fixed prices.

Let $I$ be the active link-color coordinates and $J$ the row-color
coordinates, with unselected rows set to zero. Their anchors are the coarse parent of the link
and the anchor of the row, respectively. Write
\[
             M:\mathbb R^I\longrightarrow\mathbb R^J,
             \qquad Q=M^*M
\]
for the actual selected derivative matrix and its Gram matrix, including
the fixed row normalization $\lambda>0$. Thus $M$ is the matrix obtained
by applying the selected row derivative to each active coordinate
impulse, not a new linearization at the identity. The precision completed
by the inactive identity is the matrix of Section~\ref{sec:sector-gram};
$Q$ is its restriction to active coordinates.
For matrices between either pair of anchor sets, use the weighted maximum
of absolute row and column sums, denoted $\|\cdot\|_v$ as in
Section~\ref{sec:local-covariance}. For each $j\in I$, let $W_j$ be its
actual panel. Let $N$ be the matrix whose $j$th column is the inverse
of $Q_{W_j}$, extended by zero, and let $C=(N+N^*)/2$.

\begin{proposition}[Uniform global prices]
\label{prop:actual-global-row-prices}
There are $\eta>0$, $K\ge1$, $u,v>0$, and $L_0>0$, depending only
on the fixed data, such that eventually in $x$, uniformly over the
choices just specified,
\begin{gather}
 Q\ge\eta I,\qquad C\ge\eta I,\qquad
 \|Q\|_0,\ \|C\|_0,\ \|N\|_0\le K,
 \label{eq:global-row-kernels}\\
 \|N-Q^{-1}\|_0\le u e^{-v\lfloor B/60\rfloor},\qquad
 \|M\|_{\sigma/2}\le L_0,
 \label{eq:global-row-error}
\end{gather}
where $\sigma>0$ is the fixed path width. In particular $Q$ and $C$
are positive definite. Given any fixed $r\ge0$, there is also $Z>0$
such that, for $0<d\le rt$, the actual source vector $s(Y)\in\mathbb R^J$
created by an arbitrary frozen field $Y$ obeys
\begin{equation}\label{eq:global-row-source}
                      |s(Y)_i|\le Zt\qquad(i\in J).
\end{equation}
The same tail in $x$ may be used for all these inequalities. The
source constant and tail may additionally depend on $r$.
\end{proposition}

\subsection{The actual matrix and its principal windows}

The proof of the weighted estimate in Section~\ref{sec:sector-gram}
applies to every selected row and active coordinate. Its inputs are the
uniform derivative estimate for the predictor on its actual branch and
the exponentially summable path rows. Both are valid at the coarse field
$V$ because row selection clears the required neighbourhood of its flags.
Reversing a bond changes its parent anchor by at most one; the proof
there includes this displacement. Restrictions to selected rows or active
coordinates delete nonnegative terms from the absolute sums.

More explicitly, the entry estimate there has the form
\[
 e^{\sigma\operatorname{dist}(a_i,a_j)}|M_{ij}|\le C_0.
\]
Let $m_I,m_J$ bound the numbers of input and output coordinates per
coarse anchor. They are finite constants fixed after $L$, including
all color, orientation and block-residue factors. For
$\mathcal W(v)=\sum_{n\ge0}(2n+1)^4e^{-vn}$, grouping by anchor
distance gives
\[
 \|M\|_{\sigma/2}
 \le C_0\max\{m_I,m_J\}\mathcal W(\sigma/2)=:L_0.
\]
Increasing $L_0$ if needed makes it positive. This is a global row
and column bound; its proof does not introduce the torus cardinality.
The same product estimate gives $\|Q\|_{\sigma/2}\le L_0^2$.
The coercivity argument in Section~\ref{sec:sector-gram} gives
$Q\ge\eta_0 I$ for some $\eta_0>0$. Indeed extend any active vector
by zero in the completed precision matrix; the inactive identity
contributes zero, and the quadratic form is exactly $\|M z\|_2^2$.

Set $K_0\ge\max\{1,L_0^2\}$. Apply the inverse estimates of
Section~\ref{sec:local-covariance} to this actual $Q$ and its principal
submatrices. They give $w>0$ and $S>0$, independent of the period,
for which
\[
                \|Q^{-1}\|_w\le S,\qquad \|N\|_w\le S.
\]
For the latter assertion the windows vary by column, but every
principal inverse has the same entrywise exponential bound; summing
that bound over either anchor set proves both sums. This is why a
bound on each column's Euclidean norm alone would be insufficient.
We may decrease $w$ to ensure $w\le\sigma/2$.

Each actual panel contains every coordinate within anchor distance
$\rho=\lfloor B/60\rfloor$ of its center, once $B\ge120$.
Indeed the actual panel requires every pair of coarse bond endpoints
to be at distance at most $\lfloor B/20\rfloor$. Each endpoint is at
distance at most one from its anchor, and
$2+\lfloor B/60\rfloor\le\lfloor B/20\rfloor$ for $B\ge120$.
The triangle inequality proves the asserted containment. Equivalently one
may use any family of panels with this ball-containment property.
The principal-inverse equation implies
\[
            E:=QN-I,\qquad E_{ij}=0
                 \quad\hbox{if }\operatorname{dist}(a_i,a_j)\le\rho.
\]
The product norm gives $\|E\|_w\le K_0S+1$. Remove its exponential
weight only on the remaining entries to obtain
\[
 \|E\|_0\le(K_0S+1)e^{-w\rho},\qquad
 \|N-Q^{-1}\|_0\le S(K_0S+1)e^{-w\rho}.
\]
The second inequality follows from the ordered identity
$N-Q^{-1}=Q^{-1}E$.

Since $Q^{-1}$ is symmetric, averaging $N$ with its transpose does
not increase this error. Also $Q\le K_0 I$, so
$Q^{-1}\ge K_0^{-1}I$. Because $\rho\to\infty$, enlarge the
$x$ threshold until $S(K_0S+1)e^{-w\rho}\le(2K_0)^{-1}$.
Then $C\ge(2K_0)^{-1}I$. Take
\[
 \eta=\min\{\eta_0,(2K_0)^{-1}\},\quad
 K\ge\max\{1,K_0,S\},\quad
 u=S(K_0S+1)+1,\quad v=w.
\]
These choices prove \eqref{eq:global-row-kernels} and
\eqref{eq:global-row-error}. The argument works on an empty active
set as well, with all matrix inequalities interpreted on the zero
space. No local mean or source bound was used to prove these matrix
estimates.

\subsection{Selected background rows}

The first subsection of Section~\ref{sec:warning-safety} proves the
needed bound on the affine base strokes by fixed-radius endpoint gauges,
absolute path sums, the global guide bound and the local predictor
estimate for the original tags. We check its application to the actual
selected rows here. This uses only that subsection's base-stroke argument,
not its subsequent mean calculation.

A selected row has no coarse flags in its required surrounding ball.
For large $B$, this contains every fixed chart box needed by a predictor
read or by a plaquette reached by a retained path. The actual flag
identity in the proposition therefore gives the local small plaquette
condition used by the axial gauge. Denote the fixed box constant by
$c_*$. Its coarse gauge and the parent-lifted replica gauge obey
\begin{equation}\label{eq:global-selected-gauge}
 r(V^g_e)\le c_*\operatorname{temp}(x),\qquad
 r((V^\uparrow)^{g\circ\mathrm{parent}}_e)
                          \le c_*\operatorname{temp}(x).
\end{equation}
A replica link is either the identity or a coarse link. In a signed
channel, pull back the gauge; reversal moves an anchor by at most one
and preserves chord under inversion. The extra layer in each fixed box
pays this displacement. The period condition makes the boxes nonwrapping.

The cited argument may consequently be applied separately at each path
target, since only color norms are compared. Four-link curvature costs
at most $4c_*\operatorname{temp}(x)\le8c_*t$; its absolute path sum
has a fixed price. The predictor uses its fixed local logarithm chart,
with size $O(t)$, and the same original tags. The subtracted fine and
predictor guides are precisely those of \eqref{eq:guide-maximum};
their global bounds do not require flags to vanish at an unselected
site. Thus the cited proof gives, for the unscaled affine base row,
\begin{equation}\label{eq:global-selected-base}
                      \|b_i(V)\|\le C_b t
                      \quad\hbox{on every selected row }i,
\end{equation}
with $C_b>0$ fixed before the stores, period and field. Selection uses
no assumption about the stores' fine costs or their other noise entries.
This identifies the actual base rows to which the canonical argument
applies without equating them to their flat approximations.

\subsection{Directions that include frozen links}

The source contains the derivative evaluated on the frozen displacement,
which need not be supported on the active links. We require a separate
supremum estimate on all such directions. Let $a$ be any fine-link
Lie-algebra field with $\sup_e\|a_e\|\le A$, $A\ge0$.
At a plaquette, the directional derivative is a sum of four transported
link directions, followed by the fixed curvature projection. Hence its
norm is bounded by $C_f A$, uniformly in the base configuration.
One can equivalently sum the norms over the fixed neighbouring fine
bonds; their number is bounded in terms of $L$.

For a transported path row, the differentiated curvature inputs are
summed with its absolutely summable path coefficients and isometric
adjoint transports. Thus its bound is the coefficient row sum times
$C_f A$. The direct row is bounded by its fixed square-root coefficient
times $C_f A$. These are the two components of the selected linear
curvature row. For the predictor derivative, use the local branch
estimate at \eqref{eq:global-selected-gauge}. It is bounded by a
constant times the sum of $\|a_e\|$ on the fixed predictor stencil,
hence by $C_p A$. Conjugating directions by the local gauge or applying
a signed channel preserves their norms. The same constants therefore
apply to the original direction, including its values on frozen links.
There is a $C_a>0$ such that
\begin{equation}\label{eq:global-selected-bump}
       \|\mathcal D_i(V)a\|\le C_a A
                         \quad\hbox{on every selected row }i,
\end{equation}
where $\mathcal D_i(V)$ is the derivative used in the actual row
construction. Its path transports are held as prescribed in that
construction; the estimate does not differentiate a newly chosen gauge.

Finally the frozen input $a(Y)$ is the patched displacement on frozen
links and zero on the other links. Its clipping gives
$\|a(Y)_e\|\le2d\le2rt$ for every $Y$ and every link.
The exact source is the selected restriction of
$\lambda\{b(V)+\mathcal D(V)a(Y)\}$ in scalar color coordinates.
Projection onto a color coordinate has norm at most one. Equations
\eqref{eq:global-selected-base} and \eqref{eq:global-selected-bump}
therefore give
\[
            |s(Y)_i|\le\lambda(C_b+2rC_a)t.
\]
Taking $Z=\lambda(C_b+2rC_a)+1$ proves
\eqref{eq:global-row-source}. Intersect the finitely many eventual
tails used for the matrix, base, branch and source bounds. All their
constants were fixed before the period, stores, field, tags and frozen
input, proving the asserted uniformity. No differentiability of the
frozen field as a function of $V$ was used or concluded.
